% BOOK-FILE-ID: f-0497b0842775
\subsection{特性方程式の解が虚数の場合}\label{節：特性方程式の解が虚数の場合}

\sectionref{節：隣接3項型}では,特性方程式の2解$\alpha,\beta$を使って
\[
  a_n=A\alpha^{n-1}+B\beta^{n-1}
\]
と表した.
では,特性方程式の解が虚数になったらどうなるだろうか.
実は,この式は2解が虚数でもそのまま成り立つ.
初めの2項も漸化式の係数も実数なら,各項は実数のはずである.
しかし,複素数の累乗が並んだ式からは,実数であることが分かりにくい.
そこで,まずこれまでと同じ方法で係数を決め,三角関数を使って実数の式に直してみよう.\footnote{複素数の表示や極形式は\sectionref*{節：複素平面と極形式}で確認できる.}

以下では
\[
  a_{n+2}+pa_{n+1}+qa_n=0\quad(n\ge1)
\]
の$p,q,a_1,a_2$を実数とし,$p^2-4q<0$の場合を考える.
\bookterm{characteristic}{特性方程式}$x^2+px+q=0$の2解は,互いに\bookterm{conjugate}{共役}な虚数$\alpha,\overline\alpha$となる.

\subsubsection{係数を決めてから実数の式に直す}
\[
  a_n=A\alpha^{n-1}+B\overline\alpha^{\,n-1}
\]
に$n=1,2$を代入すると,
\[
  A+B=a_1,\qquad A\alpha+B\overline\alpha=a_2
\]
である.この連立方程式を解けば,
\[
  A=\frac{a_2-\overline\alpha a_1}{\alpha-\overline\alpha},
  \qquad
  B=\frac{\alpha a_1-a_2}{\alpha-\overline\alpha}
\]
と決まる.$\alpha\ne\overline\alpha$なので割ってよい.
初期値が実数だから,この2つの係数も互いに共役になっている.

例えば$a_1=0,a_2=1$, $a_{n+2}=2a_{n+1}-2a_n$では,2解は$1\pm i$である.
$A+B=0$, $(1+i)A+(1-i)B=1$より,
\[
  A=-\frac{i}{2},\qquad B=\frac{i}{2}
\]
なので,
\[
  a_n=-\frac{i}{2}(1+i)^{n-1}+\frac{i}{2}(1-i)^{n-1}
\]
を得る.ここまででも一般項は求まっている.
次に,$1\pm i$を\bookterm{polar}{極形式}で表すと,
\[
  1\pm i=\sqrt2\left(\cos\frac{\pi}{4}\pm i\sin\frac{\pi}{4}\right)
\]
である.ド・モアブルの定理を使い,$t=\frac{(n-1)\pi}{4}$と置けば,
\[
  \begin{aligned}
    a_n&=(\sqrt2)^{n-1}\left\{-\frac{i}{2}(\cos t+i\sin t)\right.\\
    &\qquad\left.+\frac{i}{2}(\cos t-i\sin t)\right\}\\
    &=(\sqrt2)^{n-1}\sin t
  \end{aligned}
\]
となり,実数だけの式に直せた.

一般にも,虚部が正の解を
\[
  \alpha=\rho(\cos\theta+i\sin\theta),
  \qquad \rho>0,\quad 0<\theta<\pi
\]
と表す.$\rho$は原点からの距離,$\theta$は\bookterm{argument}{偏角}である.
共役な係数を,実数$C,D$を用いて
\[
  A=\frac{C-iD}{2},\qquad B=\frac{C+iD}{2}
\]
と書けば,$A+B=C$, $i(A-B)=D$なので,
\[
  \begin{aligned}
    a_n=\rho^{n-1}\{&C\cos((n-1)\theta)\\
      &+D\sin((n-1)\theta)\}
  \end{aligned}
\]
となる.共役な2項の虚部が打ち消し合い,余弦と正弦が残ったのだ.
\BookFigComplexRoots

\subsubsection{漸化式を変形して三角関数を見つける}\label{解釈：漸化式を変形して三角関数を見つける}
今度は,係数$A,B$を求めずに,元の漸化式から三角関数を見つけてみよう.
解と係数の関係より,
\[
  p=-2\rho\cos\theta,\qquad q=\rho^2
\]
だから,元の式は
\[
  a_{n+2}-2\rho\cos\theta\,a_{n+1}+\rho^2a_n=0
\]
と書ける.ここで$b_n=\frac{a_n}{\rho^{n-1}}$と置き,両辺を$\rho^{n+1}$で割ると,
\[
  b_{n+2}+b_n=2\cos\theta\,b_{n+1}
\]
となる.これは\bookterm{addition}{加法定理}
\[
  \begin{aligned}
    \cos(x+\theta)+\cos(x-\theta)&=2\cos\theta\cos x,\\
    \sin(x+\theta)+\sin(x-\theta)&=2\cos\theta\sin x
  \end{aligned}
\]
と同じ形である.$x=n\theta$を代入すれば,$\cos((n-1)\theta)$も$\sin((n-1)\theta)$もこの漸化式を満たす.
したがって,
\[
  h_n=C\cos((n-1)\theta)+D\sin((n-1)\theta)
\]
も解になる.初めの2項を$b_n$と合わせると,
\[
  C=a_1,\qquad
  D=\cfrac{\cfrac{a_2}{\rho}-a_1\cos\theta}{\sin\theta}
\]
と決まる.$0<\theta<\pi$より$\sin\theta\ne0$なので割ってよい.
$b_1=h_1,b_2=h_2$であり,$b_n=h_n,b_{n+1}=h_{n+1}$なら,
\[
  b_{n+2}=2\cos\theta\,h_{n+1}-h_n=h_{n+2}
\]
となる.帰納法により全項が一致し,$a_n=\rho^{n-1}b_n$に戻せば先ほどと同じ一般項を得る.

最初の例なら$\rho=\sqrt2$, $\theta=\frac{\pi}{4}$である.
$a_1=0$から$C=0$, $a_2=1$から$D=1$となり,
\[
  a_n=(\sqrt2)^{n-1}\sin\frac{(n-1)\pi}{4}
\]
が得られる.$A,B$を先に決める方法と,三角関数の形にしてから$C,D$を決める方法があるのである.

\begin{bookpoint}
  \textbf{POINT}\par
  特性方程式の虚数解を$\rho(\cos\theta\pm i\sin\theta)$と極形式で表し,
  $a_n=\rho^{n-1}\{C\cos((n-1)\theta)+D\sin((n-1)\theta)\}$と置いて初期条件から$C,D$を決める.
\end{bookpoint}

\subsubsection{一般項を予想して帰納法で示す}
一般項を容易に予想できる場合には,特性方程式を解く前に,その予想を帰納法で確かめることもできる.
例えば$a_1=a_2=1$, $a_{n+2}=a_{n+1}-a_n$の初めの並びは
\[
  1,1,0,-1,-1,0,1,1,\ldots
\]
である.正弦の値と比べると,$a_n=\frac{2}{\sqrt3}\sin\frac{n\pi}{3}$と予想できる.
初めの2項では確かに成り立つ.$n=k,k+1$で成り立つと仮定すると,加法定理より,
\[
  \begin{aligned}
    a_{k+2}
    &=\frac{2}{\sqrt3}\left(\sin\frac{(k+1)\pi}{3}-\sin\frac{k\pi}{3}\right)\\
    &=\frac{2}{\sqrt3}\sin\frac{(k+2)\pi}{3}
  \end{aligned}
\]
となる.したがって,数学的帰納法によりすべての項で成り立つ.
帰納法自体は一般にも使えるが,この方法で候補を見つけやすいのは,短い周期などが見える場合である.
例題(1)でも,初めの並びから一般項を予想してみよう.

このように現れる振動は,連立型の複素数の重みや,一次分数型の回転ともつながっている.
その関係と振幅・周期の詳しい説明は,付録の\sectionref*{節：虚数解型の別解と解釈}に託す.

\noindent \bookblocklink{example-complex-1}{problem}{answer}{【例題】}\par\nobreak
\begin{enumerate}[label=(\arabic*),parsep=3pt,format=\bookitemlink{example-complex-1}{problem}{answer}]
  \item $a_1=1,\,a_2=0,\quad a_{n+2}=-a_n\quad(n\ge1)$
  \item \BookMathLayout{$a_1=0,\,a_2=1,\quad a_{n+2}=2a_{n+1}-2a_n\quad(n\ge1)$}{$a_1=0,\,a_2=1$\\ $a_{n+2}=2a_{n+1}-2a_n$\\ $(n\ge1)$}
  \item $a_1=1,\,a_2=2,\quad a_{n+2}=-4a_n\quad(n\ge1)$
  \item \BookMathLayout{$a_1=1,\,a_2=0,\quad a_{n+2}=2a_{n+1}-4a_n\quad(n\ge1)$}{$a_1=1,\,a_2=0$\\ $a_{n+2}=2a_{n+1}-4a_n$\\ $(n\ge1)$}
\end{enumerate}

\noindent\bookblocklink{example-complex-1}{hint}{problem}{【方針】}\par\nobreak
\begin{enumerate}[label=(\arabic*),parsep=3pt,format=\bookitemlink{example-complex-1}{hint}{problem}]
  \item 初項から数項を計算して\bookterm{periodic}{周期性}を見つけ,\bookterm{general}{一般項}を予想します.その後,\bookterm{induction}{数学的帰納法}で証明します.
  \item 特性方程式の解を\bookterm{polar}{極形式}で表します.$\rho=\sqrt2$,$\theta=\frac{\pi}{4}$として\bookterm{general}{一般項}の形を置きます.
  \item 特性方程式の解は$\pm2i$です.$\rho=2$,$\theta=\frac{\pi}{2}$として\bookterm{initial}{初期条件}を代入します.
  \item 特性方程式の解を極形式で表し,$\rho=2$,$\theta=\frac{\pi}{3}$として係数を決定します.
\end{enumerate}

\noindent\bookblocklink{example-complex-1}{answer}{problem}{【解答】}\par\nobreak
\begin{enumerate}[label=(\arabic*),parsep=3pt,format=\bookitemlink{example-complex-1}{answer}{problem}]
  \item $a_n=\cos\frac{(n-1)\pi}{2}$
  \item $a_n=(\sqrt2)^{n-1}\sin\frac{(n-1)\pi}{4}$
  \item $a_n=2^{n-1}\left(\cos\frac{(n-1)\pi}{2}+\sin\frac{(n-1)\pi}{2}\right)$
  \item \BookMathLayout{$a_n=2^{n-1}\left(\cos\frac{(n-1)\pi}{3}-\frac{1}{\sqrt3}\sin\frac{(n-1)\pi}{3}\right)$}{$a_n=2^{n-1}\Bigl(\cos\frac{(n-1)\pi}{3}$\\
      $\hspace{3zw}-\frac{1}{\sqrt3}\sin\frac{(n-1)\pi}{3}\Bigr)$}
\end{enumerate}

\noindent\bookblocklink{example-complex-1}{solution}{problem}{【解法】}\par\nobreak
\begin{enumerate}[label=(\arabic*),parsep=3pt,format=\bookitemlink{example-complex-1}{solution}{problem}]
  \item 初めの並びは$1,0,-1,0,1,0,-1,0,\ldots$なので,
    $a_n=\cos\frac{(n-1)\pi}{2}$と予想する.
    $n=1,2$で初期値に合い,$a_k=\cos\frac{(k-1)\pi}{2}$なら,
    \[
      \begin{aligned}
        a_{k+2}&=-a_k=\cos\left(\frac{(k-1)\pi}{2}+\pi\right)\\
        &=\cos\frac{(k+1)\pi}{2}.
      \end{aligned}
    \]
    従って,奇数番目と偶数番目それぞれの\bookterm{induction}{数学的帰納法}で全項について成り立つ.
  \item \bookterm{characteristic}{特性方程式}$x^2-2x+2=0$の根は$1\pm i$.
    $\rho=\sqrt2,\theta=\frac{\pi}{4}$である.
    $a_1=0$から$C=0$,$a_2=1$から$D=1$と決まり,【解答】を得る.
  \item $x^2+4=0$の根は$\pm2i$なので,$\rho=2,\theta=\frac{\pi}{2}$.
    $a_1=1$から$C=1$,$a_2=2$から$2D=2$,従って$D=1$となる.
  \item $x^2-2x+4=0$の根は$1\pm\sqrt3\,i$なので,$\rho=2,\theta=\frac{\pi}{3}$.
    $a_1=1$から$C=1$,$a_2=0$から
    \[
      0=2\left(\frac{1}{2}+D\frac{\sqrt3}{2}\right)
    \]
    となり,$D=-\frac{1}{\sqrt3}$を得る.
\end{enumerate}
